\subsubsection{The fitted starting fields for the even-order case}
\label{par:count-ii-tower}

To apply propagation in even order, the starting field must retain enough marked Galois elements to separate the ordinary labels. This requirement governs which character kernels one may impose.

\begin{proposition}[Starting field and marked groups]\label{prop:count-ii-starting-field}
Assume the even-order case of
Lemma~\ref{lem:odd-prime-positive-ranges} and use its auxiliary
character from Lemma~\ref{lem:marked-dyadic-character}.
Write \(H_0\) for the subgroup cut out by the finite character conditions in \eqref{eq:count-ii-starting-groups}. The starting fields are specified by 
\[
 G_{F_0}=H_0,\qquad G_K=\ker\lambda,\qquad G_{k_0}=\ker\phi_0.
\]
 The marked normalizing and inertia elements
are retained, and the chosen quadratic extension splits at the
coefficient places used for the ordinary flag.
\end{proposition}

\begin{proof}
\proofstep{The rational characters and the starting field.}

Use the rational character data of
Paragraph~\ref{par:dyadic-residual} and the \(\eta\) of
conductor \(2^n\) there. The order of the prime-to-\(\ell\)
part of \(\xi\) away from \(\ell\) is even in this case.
Write \(\kappa_0\) for the \(\ell\)-primary part of \(\xi\)
away from \(\ell\), and \(\ell^b\) for an exponent killing
it. These factors are taken first on the character group
and then on the two coprime conductor factors. In
particular \(\kappa_0\) is unramified at \(\ell\) and away
from \(S\).

Let \(m_0\) be the order of the prime-to-\(\ell\)
part of \(\xi\). Write \(\xi=\xi_{\ell'}\kappa_0\xi_{\ell,\ell}\), where the last factor is the \(\ell\)-primary component at \(\ell\). On the simultaneous kernel we have
\[
 \kappa_0(g)=1,\qquad \epsilon_\ell(g)=1\Longrightarrow\xi_{\ell,\ell}(g)=1,
 \qquad \xi(g)^{m_0}=\xi_{\ell'}(g)^{m_0}=1 .
\] On inertia at a place away from
\(\ell\), the same implication holds using just
\(\kappa_0=1\). On that part of inertia the residual cyclotomic factor is
trivial, so \(\bar\Psi=1\) implies \(\bar\xi=1\). Reduction is
injective on prime-to-\(\ell\) roots, giving
\[
              \xi^{m_0}=1,\ \bar\xi=1\ \Longrightarrow\ \xi=1 .
\]

Choose three quadratic signs \(\phi_0,\lambda,\lambda'\)
of distinct odd prime conductors \(p_3,p,p'\) away
from \(S\), with parities \(+,-,-\), respectively.
To do so, given a parity \(s\in\{1,-1\}\) and
the finite set of primes \(S'\) already used or
required to split, take a prime
\(P\equiv s\pmod{8\prod_{q\in S'}q}\) outside \(S'\)
by Dirichlet's theorem. For odd \(q\in S'\), quadratic reciprocity gives
\[
 \left(\frac qP\right)
 =\left(\frac{s}{q}\right)(-1)^{((q-1)/2)((P-1)/2)}=1 .
\]

For \(s=-1\) the two signs are each \((-1)^{(q-1)/2}\), and for \(s=1\) each is positive. Also \((2/P)=1\) since \(P\equiv\pm1\pmod8\). The quadratic character
of conductor \(P\) has parity \(s\), is ramified at
\(P\), and is unramified elsewhere at finite places.
Start with \(S'=S\) and add each chosen prime before
the next choice. An inertia element at any of the
three conductors can change the corresponding sign
independently. It fixes the values of \(\xi,\Psi,\omega,\eta,\kappa_0\),
since these are unramified at the three primes.

In \(G_{\Q}\) take the open subgroup
\begin{equation}\label{eq:count-ii-starting-groups}
 H_0=\{g:\bar\Psi(g)=\phi_0(g),\quad
             \lambda(g)=\lambda'(g),\quad\kappa_0(g)=1\}.
\end{equation}
The signs in the first equality are taken in the
residue field. This subgroup is open because the
three conditions are finite character equalities,
and it contains the commutator subgroup. It contains
complex conjugation, since \(\Psi\) is even,
\(\lambda,\lambda'\) are odd and \(\ell\) is odd.
Its index is even since \(\lambda/\lambda'\) has
order two.

Let \(F_0\), \(K\) and \(k_0\) be the fixed
fields of \(H_0\), \(\ker\lambda\) and \(\ker\phi_0\),
respectively. The fixed fields have
\[
 \begin{gathered}
 F_0/\Q\text{ abelian and totally real},\qquad 2\mid[F_0:\Q],\\
 K/\Q\text{ imaginary quadratic},\qquad
 k_0/\Q\text{ real quadratic}.
 \end{gathered}
\]
The chosen congruences make \(K\) split at \(S\) and \(k_0\)
split at \(2,\ell\). To check the ramification of \(L_0=F_0K\) over \(F_0\),
for a rational inertia group \(I_v\) use \(I_v\cap H_0\).
At \(p\), \(\lambda'|_{I_p}=1\) and
\(\lambda=\lambda'\) on \(H_0\); away from \(p\), \(\lambda\)
is already unramified. Thus
\[
                   \lambda|_{I_v\cap H_0}=1
                         \qquad(v\text{ finite}).
\]

At \(p_3\) use \(\bar\Psi|_{I_{p_3}}=1\) because \(p_3\) lies
outside the host conductor. The other places are already unramified
for \(\phi_0\). Therefore
\[
 \phi_0|_{I_v\cap H_0}=1\quad(v\text{ finite}),
\]
so both quadratic composita are unramified over \(F_0\) there. The field \(F_0\) is tamely
ramified at \(\ell\), because on inertia there all
three quadratic signs and \(\kappa_0\) are trivial,
and the residual \(\bar\Psi\) has order prime to
\(\ell\). Inertia of \(F_0\) at \(\ell\) kills the
entire diagonal residual representation since
\(\eta\) is unramified there. Above each prime of
\(F_0\) dividing \(\ell\), choose its extension to
\(L_0\) inducing the original marked place of \(K\).
Exactly one of the two extensions does so, since
\(K\) splits at \(\ell\). Choose decomposition
embeddings representing these places in the fixed
algebraic closure.

The subgroup contains local elements with the
required large values of \(\bar\eta\). At \(\ell\) take
the Frobenius lift \(f_\ell\) with \(\epsilon_\ell(f_\ell)=1\).
Choose \(2^a u\), \(u\) odd, killing \(\xi\), as in
the dyadic calculation. With
\[
 e=2^a u\ell^b,
\]
the chosen lift satisfies
\[
 \epsilon_\ell(f_\ell^e)=1,\qquad
 \xi(f_\ell^e)=1,\qquad \kappa_0(f_\ell^e)=1.
\]
It therefore kills \(\Psi\), and all three auxiliary signs are trivial on this decomposition group. If the \(\bar\eta\)-value had order at most four, injectivity
of reduction would give \(\eta(\ell)^{2^{a+2}u\ell^b}=1\).
The dyadic test then gives
\[
                 2^n<\ell^2 2^{a+2}\le 2^{n_{\rm II}},
\]
contrary to \(n=n_{\rm II}\).

At two use inertia with cyclotomic
residue \(3\bmod 2^n\) for the \(2^n\)-th roots
of unity. It exists by total ramification there.
Raise it to \(2^{c_2(\xi)}\ell^b\) to kill the
host and \(\kappa_0\). If its \(\bar\eta\)-value had order at most four, the same
test with residue \(3\) would give
\[
       2^n<9\cdot2^{c_2(\xi)+2}<2^{c_2(\xi)+6}\le 2^n .
\]
Thus its order is greater than four.
Its \(\omega\)-value is one. Denote this element by
\(u_0\). Both local elements belong to \(G_{L_0}\).
They are available at every place in question by
conjugation, since the groups defining the fields
are normal and the characters are abelian.

We need two further elements of \(G_{L_0}\) with prescribed character values.
Complex conjugation has \(\xi,\Psi,\eta,\omega\)-values
\((\xi(-1),1,-1,-1)\). Set the three auxiliary signs
to \((1,1,1)\) by the inertia adjustments above.
The resulting element \(a_0\) has
\[
 (\bar\eta,\bar\Psi,\bar\omega)=(-1,1,-1).
\]

To construct the second element, recall that the even-order prime-to-\(\ell\)
factor of \(\xi\) off \(\ell\) takes the value \(-1\).
Choose an element with this value and
\(\epsilon_\ell=1\), using the independent cyclotomic
quotients at coprime moduli. Raise it to \(\ell^b\)
to kill \(\kappa_0\) and set the auxiliary signs to
\((-1,1,1)\). The adjustments preserve
\(\epsilon_\ell=1\). If the resulting element \(g\) has \(\bar\eta(g)^2=1\), replace it by \(gu_0\). Indeed
\[
 \bar\eta(gu_0)^2=\bar\eta(u_0)^2\ne1,\qquad
 \bar\Psi(gu_0)=\bar\Psi(g),\qquad \epsilon_\ell(gu_0)=1 ,
\]
using the marked values of \(u_0\) and triviality of the cyclotomic character on inertia at two.

We obtain the marked normalizing element \(g_0\)
with \(\bar\Psi(g_0)=-1\), \(\epsilon_\ell(g_0)=1\) and
\(\bar\eta(g_0)^2\ne1\). The labels on it are the
four distinct values \(-\bar\eta(g_0),-1,\bar\eta(g_0),1\),
in the order allowed by the rational chart. In
particular the quadratic sign stays nontrivial on
\(G_{L_0}\cap\ker\epsilon_\ell\).

The values needed for separation are
\[
 \begin{array}{c|ccc}
       &\bar\eta &\bar\Psi &\bar\omega\\ \hline
 a_0   &-1 &1 &-1\\
 u_0   &\zeta &1 &1\\
 g_0   &z &-1 &1
 \end{array}
 \qquad \operatorname{ord}(\zeta)>4,\quad z^2\ne1.
\]
The order of \(\zeta\) is a power of two. For even \(i\)
the torus character \(\bar\eta^i\bar\Psi^j\) is \(1\) at
\(a_0\), whereas for odd \(i\) it is \(\zeta^i\ne1\) at
\(u_0\). Both cases distinguish the character from \(\bar\omega\). Every root character
is nontrivial, using \(u_0\) for odd first exponent
and \(g_0\) for the two roots \(\pm w_\Psi\).
The first root also has order greater than four.
Suppose a residual quadratic character trivial on
\(\ker\omega\cap G_{L_0}\) permutes the four residual
labels by multiplication. Applied to the trivial
label it must equal one of the labels.

Since \(u_0,g_0\in\ker\omega\), a putative twisting label \(\lambda\) must satisfy
\[
 (\lambda(u_0),\lambda(g_0))=(1,1).
\]
Both labels involving \(\bar\eta\) fail the first equality and the other nontrivial label fails the second. Thus \(\lambda=1\). At each place above \(\ell\) the local
test of nontrivial \(\bar\eta\) with trivial sign and
Teichm\"uller value is given by the Frobenius power
constructed above. The quadratic restriction also
makes \(\eta|_{G_{k_0}}\) totally odd and unramified
at the two degree-one places above \(\ell\). Its
value on their Frobenius has order greater than
four, as does its restriction on inertia at either
norm-two place.

\end{proof}

\subsubsection{Auxiliary inertia and deeper layers}

The starting field supplies the marked elements. Enlarging it to kill auxiliary inertia must still leave enough primes with the same residual Frobenius behavior for the two level comparisons.

\begin{proposition}[Auxiliary primes on the deep layer]\label{prop:count-ii-auxiliary-primes}
After enlarging the even-order starting field to kill the added
inertia, one can choose arbitrarily many distinct rational primes \(p\), outside the fixed bad set, splitting in \(L\). For a Frobenius \(f_p\) the two auxiliary levels use the simultaneous conditions 
\[
 \bar r(f_p)\text{ is scalar},\qquad \bar\epsilon_\ell(f_p)\ne1.
\]
 Both levels use the same scalar twist of \(\bar r\); twisting multiplies the Frobenius matrix by a scalar.
\end{proposition}

\begin{proof}
The marked cup construction of
Paragraph~\ref{par:marked-cup} enlarges \(F_0\) by
the degree-\(\ell\) cyclotomic fields of the added
rational primes. Write \(F_1\) for the resulting
compositum. This is a finite abelian extension of
odd degree over \(F_0\), unramified at \(\ell\).
It kills the new inertia in the residual chain
and preserves the tame property of the base
field at \(\ell\). In every subsequent abelian
\(\ell\)-power extension of \(F_1\), a common
\(\ell\)-power of the marked elements belongs to
the deeper absolute Galois group. Every marked finite value \(z\) used for these comparisons
has order prime to \(\ell\), so
\[
             \operatorname{ord}(z^{\ell^a})=\operatorname{ord}(z)
\]
for all \(a\). The powers retain the required distinctions.

Taking these powers also preserves the specified decomposition and inertia subgroups. The exact
conditions \(\epsilon_\ell=1\) and \(\omega=1\)
on the indicated elements persist too. The local
frames are transported using the selected places
on the same side of \(K\). The normalizing element
at a new layer is the corresponding power of
\(g_0\), so the residual labels stay distinct there.

Use the cohomological bound \(C\) chosen on \(F_1\)
as in the component argument, before extending
the field further. To make the host restriction
have order prime to \(\ell\), take the compositum
of \(F_1\) with the degree-\(\ell^m\) layer of the
cyclotomic \(\mathbf Z_\ell\)-extension of \(\Q\).
Its absolute Galois subgroup is
\[
 G_{F_1}\cap\{g:\epsilon_\ell(g)^{\ell-1}
                        \equiv1\pmod{\ell^{m+1}}\}.
\]
The kernel of \(\epsilon_\ell^{\ell-1}\) on
\(G_{F_1}\) lies in \(\ker\xi^{m_0(\ell-1)}\).
Indeed \(g^{\ell-1}\) has cyclotomic value one
there and \(G_{F_1}\subseteq\ker\kappa_0\). This
latter kernel is open. Writing \(H_m\) for the displayed nested subgroups, their
intersection is \(\ker(\epsilon_\ell^{\ell-1}|_{G_{F_1}})\).
They are compact, so the finite character kernel gives
\[
         H_m\subseteq\ker\xi^{m_0(\ell-1)}
                  \qquad\text{for all sufficiently large }m .
\]
Otherwise the nonempty closed sets of counterexamples would have
nonempty intersection, contrary to the assertion at the limit.

The exponent \(m_0(\ell-1)\) is prime to \(\ell\).
For degree growth choose a positive prime-to-\(\ell\)
power \(s\) of the marked inertia element of
infinite cyclotomic order such that \(s\in G_{F_1}\).
This is possible by tameness. For a prescribed \(B\), none of the finitely many \(s^k\), \(1\le k<B\), lies in the limiting subgroup, since their cyclotomic images have infinite order. The subgroups decrease, so after taking one common depth,
\[
 s^iG_F=s^jG_F\quad\Longrightarrow\quad s^{j-i}\in G_F
 \quad\Longrightarrow\quad i=j\qquad(0\le i,j<B).
\]
Choose \(B=30(C+1)^2\). Since the \(B\) powers have distinct images at all sufficiently
deep levels,
\[
 [F:\Q]\ge[G_{F_1}:G_F]\ge B.
\]
Increase the depth
once for this bound and the finite character condition; the fitted
data are retained. The seed
and the three residual lines can now be formed
at this layer with the transported labels, the
same \(K\), and the real quadratic character
\[
 \phi_0|_{G_F}.
\]


The transported \(a_0\) still has \(\bar\omega=-1\),
so \(L=FK\) does not contain the \(\ell\)-th roots
of unity. Take a power of the transported \(u_0\)
whose \(\bar\eta\)-value is \(-1\) and multiply
it by the transported \(a_0\). This gives an element
of \(G_L\) with trivial residual diagonal characters
and nontrivial mod-\(\ell\) cyclotomic value. Apply
Chebotarev to the finite quotient given by \(L\),
the residual characters and the mod-\(\ell\) cyclotomic
character. This provides arbitrarily many rational
primes split in \(L\) with these Frobenius values.
The auxiliary-prime separation used in the two
ordinary comparisons is therefore available at
distinct primes outside the fixed bad set. A scalar
twist still gives a single residual eigenvalue at
each of these primes.
\label{cor:ordinary-count-ii-tower-1}
\end{proof}
